
un homomorphisme de groupes de @@M0@@ dans @@M1@@; l'application qui à une série associe son @@M2@@-ième coefficient, multiplié par @@M3@@, est un homomorphisme additif de @@M4@@ dans @@M5@@~; le morphisme @@M6@@ sera le composé des deux. Il est clair que @@M7@@ d'après la définition (7.1.1).

La seconde des formules (7.5.1) exprime donc la commutativité du diagramme écrit plus haut. Si on suppose que @@M8@@ est un @@M9@@-anneau, le carré de gauche est commutatif, et il suffit de vérifier la commutativité de celui de droite. Or il s'agit d'un diagramme d'homomorphismes de foncteurs sur @@M10@@, ayant pour source @@M11@@~; il suffit donc d'après 1.10 de vérifier la commutativité pour @@M12@@~; on a alors :
@@M13@@
@@M14@@, ce qui donne le résultat.

Pour montrer la dernière, on remarque d'abord que le diagramme :
@@M15@@
est commutatif. Comme plus haut, il suffit en effet de le voir pour @@M16@@ et on a :
@@M17@@
@@M18@@

Par suite, @@M19@@.

Par ailleurs, @@M20@@.
